% !Mode:: "TeX:UTF-8"

\chapter{绪论}[Introduction]

\section{研究背景与意义}[Background and significance]

人形机器人由于具有接近人类的身体结构和运动形态，被认为是最有潜力在真实人类环境中执行复杂任务的机器人形态之一。与传统工业机器人主要服务于结构化生产线不同，人形机器人能够在面向人类设计的空间中完成移动、操作、交互与服务等多样化任务，例如室内服务、灾害救援、医疗辅助以及危险环境巡检等。因此，使人形机器人具备稳定、灵活且可扩展的全身运动能力，是实现通用具身智能和未来服务机器人的重要基础，典型应用场景与任务形式如图~\ref{fig:humanoid_applications_intro} 所示。

\begin{figure}[htbp]
  \centering
  \vspace{2mm}
  \includegraphics[width=0.90\textwidth]{人形机器人应用场景.png}
  \vspace{2mm}
  \caption{人形机器人典型应用场景示意图}
  \label{fig:humanoid_applications_intro}
\end{figure}

除应用需求持续扩大外，近年来国内外人形机器人本体平台也进入快速迭代阶段。英国 Engineered Arts 推出的 AMECA、美国的 Apollo、Optimus、Atlas、Figure 01，以及国内的 Fourier GR-1、宇树 G1、Walker S1、CyberOne、AgiBot 等平台，分别在交互展示、双足运动、灵巧操作和通用作业等方向展开探索。不同产品在驱动形式、整机尺寸、自由度配置、上肢能力和目标应用场景上存在明显差异，这也反映出当前人形机器人仍处于技术路线并行演进的阶段。典型国内外人形机器人产品形态如图~\ref{fig:humanoid_products_intro} 所示。

\begin{figure}[htbp]
  \centering
  \vspace{2mm}
  \includegraphics[width=0.98\textwidth]{人形机器人产品.png}
  \vspace{2mm}
  \caption{国内外代表性人形机器人产品示意图}
  \label{fig:humanoid_products_intro}
\end{figure}

过去较长时间内，人形机器人研究主要关注稳定行走、速度跟踪和基础越障等移动能力。这类研究通常将下肢视为移动平台，将控制重点放在质心稳定、足端接触和抗扰恢复上，而对躯干、上肢和全身姿态表达能力的利用相对有限。随着硬件平台、物理仿真器和学习算法的发展，研究者开始关注更具挑战性的全身敏捷运动，例如跳跃、快速转向、舞蹈、运动技能和人机交互动作等\cite{li2023robust,radosavovic2024real,cheng2024expressive,ji2024exbody2}。这些动作不再只是“保持不摔倒”的移动任务，而是要求机器人在多关节协同、多接触切换和动态平衡之间保持一致的运动风格与物理可行性。

然而，与轮式机器人或固定基座机械臂相比，人形机器人在控制上面临更为严峻的挑战。一方面，人形机器人通常具有数十个自由度，其运动控制涉及腿部支撑、躯干姿态、手臂摆动和头部朝向等多部分协同；另一方面，行走、跳跃和转身过程中会频繁发生单脚支撑、双脚支撑、腾空、落地等接触模式切换，系统呈现强非线性、强耦合和欠驱动特性。此外，地面接触、摩擦不确定性、执行器带宽限制和传感器噪声等因素会进一步增加控制难度。因此，如何在高自由度、多接触和复杂动力学约束条件下实现自然、稳定且敏捷的全身运动控制，仍是机器人领域的重要问题。

传统控制方法主要依赖精确的动力学建模、轨迹规划与反馈控制器设计，例如基于模型预测控制（Model Predictive Control, MPC）、零力矩点（Zero Moment Point, ZMP）稳定性分析以及逆动力学控制等方法。这类方法具有较好的可解释性，在结构明确、任务边界清晰的场景中能够取得稳定效果。但当任务扩展到跳跃、快速踢腿、连续转身等高动态动作时，传统方法往往需要复杂的接触序列设计、人工调参和分阶段控制逻辑。对于动作类别丰富、姿态变化剧烈的人形机器人运动技能学习而言，完全依赖手工设计的控制框架难以充分利用大规模运动数据中的先验知识。

近年来，深度强化学习（Deep Reinforcement Learning, DRL）为高维连续控制问题提供了数据驱动的解决思路。通过在物理仿真环境中进行大规模并行交互，策略网络能够直接学习从机器人状态到控制指令的映射，并在奖励函数约束下逐步形成满足动力学规律的运动模式。Isaac Gym 等图形处理器（Graphics Processing Unit, GPU）并行仿真平台显著提升了机器人强化学习的采样效率，并推动了分钟级并行训练范式的发展\cite{isaacgym2021,rudin2022learning,margolis2022rapid}，近端策略优化（Proximal Policy Optimization, PPO）等算法则为连续控制策略训练提供了相对稳定的优化框架\cite{ppo2017}。在此基础上，基于学习的方法已经在人形机器人行走、跳跃、跑酷和复杂地形运动中展现出较强潜力\cite{xie2020learning,radosavovic2402humanoid,zhuang2024humanoid,gu2024advancing,long2024learning,zhang2024whole}。

与此同时，具身智能与大模型技术的快速发展正在改变人形机器人的研究范式。近期综述指出，大语言模型、视觉 Transformer、视觉语言模型以及具身多模态大模型，正在持续增强机器人在语言理解、场景感知、常识推理和任务规划等方面的能力\cite{zeng2025embodied}。对于人形机器人而言，这意味着其控制问题不再局限于“如何稳定迈步”，而是逐步扩展为“如何理解任务语义、感知开放环境，并将高层意图转化为自然的全身动作执行”。不过，高层语义能力的提升并不能直接替代底层运动控制。面对高自由度身体结构、频繁接触切换和动态平衡约束，机器人仍然需要可靠的动作表示、可执行的参考轨迹以及稳定的闭环控制策略，才能将感知和规划结果真正落实为物理世界中的可行动作。

另一方面，人体运动数据为人形机器人学习自然运动提供了重要来源。人体表面形状动作捕捉档案（Archive of Motion Capture as Surface Shapes, AMASS）等大规模人体运动捕捉数据集将不同来源的动作捕捉（Motion Capture, MoCap）数据统一到多人线性蒙皮模型（Skinned Multi-Person Linear, SMPL）参数化人体模型中，包含丰富的根节点位姿、关节姿态和体型参数\cite{amass2019}。DeepMimic 证明了强化学习可以通过跟踪参考动作来获得高质量的物理角色技能\cite{deepmimic2018}，后续工作进一步探索了对抗运动先验、可复用技能嵌入、风格化运动控制以及基于人体参考的整身控制方法\cite{peng2021amp,peng2022ase,tessler2024maskedmimic,luo2023perpetual}。这些研究表明，将人体运动先验与物理仿真中的闭环控制结合，是获得自然、协调和高动态运动技能的有效路径。

从这一视角看，大模型技术与人体动作模仿并非割裂的两条路线，而是可以形成互补关系。大模型擅长处理开放场景中的语义理解、跨模态关联和任务分解问题，人体运动数据则为机器人提供自然、连贯且符合人类运动习惯的动作先验\cite{zeng2025embodied}。对于本文关注的敏捷全身运动技能学习任务而言，优先解决专家数据预处理、动作重定向和强化学习稳定跟踪问题，有助于为后续更高层级的语言条件控制、视觉条件控制或多模态任务控制建立可复用的底层运动能力基础。

尽管如此，直接利用人体运动数据驱动人形机器人仍面临若干关键挑战。首先，人体骨架与机器人结构在拓扑关系、连杆比例、关节轴定义和自由度数量上存在显著差异，原始人体姿态无法直接作为机器人关节目标，需要通过动作重定向（motion retargeting）转换为机器人可执行的参考轨迹\cite{he2024learning,zhang2024whole}。其次，人体运动数据通常只提供运动学信息，缺少接触力、关节力矩和执行器约束等动力学信息，若在仿真中直接跟踪，容易产生足端滑移、关节越界、扭矩异常或整体失稳等问题\cite{tessler2024maskedmimic,he2024learning}。再次，高动态动作训练过程中，策略可能通过降低动作幅度、提前终止或保持保守姿态来获得局部较高奖励，导致模仿精度和动作表现力不足。此外，策略在不同初始相位、外部扰动和环境变化下的鲁棒性仍需专门设计和评估\cite{deepmimic2018,tobin2017domain,tan2018sim}。上述关键挑战的归纳如图~\ref{fig:humanoid_learning_challenges} 所示。

\begin{figure}[htbp]
  \centering
  \includegraphics[width=0.90\textwidth]{人形机器人挑战.png}
  \caption{人体数据驱动人形机器人运动技能学习的关键挑战}
  \label{fig:humanoid_learning_challenges}
\end{figure}

针对上述问题，本文研究基于强化学习的人形机器人运动技能学习方法，重点关注人体专家数据向机器人可执行轨迹的转换、高保真物理仿真中的整身动作跟踪学习，以及策略鲁棒性与泛化能力的提升机制。本文暂不以真实机器人部署为主要目标，而是将研究重点放在仿真环境中的数据处理、参考轨迹构建、强化学习任务建模和动作跟踪效果验证上。通过构建统一的技术流程，本文旨在为人形机器人学习多样化全身运动技能提供可复现的实现框架和实验基础。

\section{研究目标与主要内容}[Objectives and contents]

本文面向人形机器人复杂运动技能学习任务，构建一套由专家数据预处理、动作重定向、强化学习动作跟踪和策略泛化评估组成的技术流程。研究目标是在高保真仿真环境中，使人形机器人能够跟踪由人体运动数据转换得到的参考动作，并在扰动和环境变化下保持较好的稳定性。

本文的主要研究内容包括以下三个方面。

第一，研究专家数据预处理与动作重定向方法。针对 AMASS 等人体运动捕捉数据，本文从 SMPL 参数化人体模型出发，解析人体根节点位姿、关节旋转与身体形状参数，建立人体关键点与目标机器人连杆关键点之间的拓扑对应关系。进一步结合体型拟合、运动学优化、关节范围约束、时序平滑以及异常片段清洗等步骤，将原始人体动作序列转换为机器人可执行、时序连续且满足基本物理可行性的参考轨迹，为后续强化学习动作跟踪任务提供高质量专家动作输入。

第二，构建基于强化学习的人形机器人整身动作跟踪方法。围绕目标机器人建立高保真仿真训练环境，统一定义状态观测、动作输出、低层比例-微分（Proportional-Derivative, PD）控制器和外部扰动注入方式，并以重定向后的参考轨迹作为专家动作输入，设计相位驱动的动作跟踪任务。在此基础上，综合构建观测空间、动作空间、参考状态表示、奖励函数以及 PPO 优化目标，使策略能够在长时序执行过程中实现对参考动作的稳定跟踪，并保持较好的动态平衡和动作自然性\cite{isaacgym2021}。

第三，研究策略稳定化与实验评估方法。针对动作跟踪过程中容易受到初始状态偏差、接触扰动和环境变化影响的问题，本文引入参考状态初始化、课程学习、终止机制和域随机化等训练稳定化策略，以提升策略的鲁棒性与泛化能力。同时，围绕动作跟踪精度、训练收敛速度、执行稳定性、抗扰能力和泛化性能等指标设计系统实验，并结合可视化结果与典型案例分析，对所提方法的有效性和适用范围进行验证。

\section{论文结构安排}[Thesis organization]

本文后续章节安排如下：第2章介绍人形机器人运动技能学习的相关研究；第3章给出专家数据预处理与动作重定向方法；第4章阐述基于强化学习的动作跟踪任务、观测与动作建模、PPO 优化目标、奖励函数和稳定化训练机制；第5章给出实验设计、评价指标与实验结果分析；第6章对全文工作进行总结，并从研究局限性与未来工作展望两个方面展开讨论。

\section{本文总体技术路线}[Overall framework]

本文总体方案包括三个阶段。

第一阶段为专家数据预处理。选取人体运动捕捉数据作为专家动作来源，提取 SMPL 运动参数，建立人体关键点与目标人形机器人连杆之间的拓扑对应关系\cite{he2024learning}。通过静态体型拟合得到与机器人几何尺寸更接近的人体表示，再通过逐帧动作重定向求解机器人关节角和根位姿轨迹。

第二阶段为强化学习动作跟踪。将重定向后的参考轨迹封装为标准化动作库，在仿真环境中实时采样参考状态。策略网络根据机器人本体状态、动作相位和参考误差输出控制动作，底层 PD 控制器将动作转换为关节控制命令。奖励函数围绕刚体位置、姿态、关节角、关节速度和安全正则项构建，训练算法采用 PPO。

第三阶段为鲁棒性与泛化测试。在训练完成后，向仿真环境中加入外部推力、传感器噪声、控制延迟和不同地形几何，评估策略在扰动条件下的动作保持能力。同时可进行跨仿真平台测试，用于检查动作跟踪方法对不同物理引擎和参数设置的敏感性。

本文基于人体数据与强化学习的人形机器人运动技能学习总体方案如图~\ref{fig:chapter1_overall_framework} 所示。

\begin{figure}[htbp]
  \centering
  \vspace{3mm}
  \includegraphics[width=0.95\textwidth]{总体方案.png}
  \vspace{3mm}
  \caption{基于人体数据与强化学习的人形机器人运动技能学习总体方案}
  \label{fig:chapter1_overall_framework}
\end{figure}
